% Options for packages loaded elsewhere
\PassOptionsToPackage{unicode}{hyperref}
\PassOptionsToPackage{hyphens}{url}
\PassOptionsToPackage{dvipsnames,svgnames,x11names}{xcolor}
\documentclass[
  10pt,
  fontset=fandol]{ctexart}
\usepackage{xcolor}
\usepackage[margin=16mm]{geometry}
\usepackage{amsmath,amssymb}
\setcounter{secnumdepth}{-\maxdimen} % remove section numbering
\usepackage{iftex}
\ifPDFTeX
  \usepackage[T1]{fontenc}
  \usepackage[utf8]{inputenc}
  \usepackage{textcomp} % provide euro and other symbols
\else % if luatex or xetex
  \usepackage{unicode-math} % this also loads fontspec
  \defaultfontfeatures{Scale=MatchLowercase}
  \defaultfontfeatures[\rmfamily]{Ligatures=TeX,Scale=1}
\fi
\usepackage{lmodern}
\ifPDFTeX\else
  % xetex/luatex font selection
\fi
% Use upquote if available, for straight quotes in verbatim environments
\IfFileExists{upquote.sty}{\usepackage{upquote}}{}
\IfFileExists{microtype.sty}{% use microtype if available
  \usepackage[]{microtype}
  \UseMicrotypeSet[protrusion]{basicmath} % disable protrusion for tt fonts
}{}
\makeatletter
\@ifundefined{KOMAClassName}{% if non-KOMA class
  \IfFileExists{parskip.sty}{%
    \usepackage{parskip}
  }{% else
    \setlength{\parindent}{0pt}
    \setlength{\parskip}{6pt plus 2pt minus 1pt}}
}{% if KOMA class
  \KOMAoptions{parskip=half}}
\makeatother
\setlength{\emergencystretch}{3em} % prevent overfull lines
\providecommand{\tightlist}{%
  \setlength{\itemsep}{0pt}\setlength{\parskip}{0pt}}
\usepackage{amsmath,amssymb,mathtools,needspace}
\xeCJKDeclareCharClass{CJK}{"2032,"0400->"04FF,"2160->"216B,"2460->"2473,"25A0,"25C7,"25CB,"25CF}
\IfFontExistsTF{Arial Unicode MS}{\xeCJKsetup{AutoFallBack=true}\setCJKfallbackfamilyfont{\CJKrmdefault}{Arial Unicode MS}}{}
\setcounter{secnumdepth}{0}
\setlength{\emergencystretch}{2em}
\allowdisplaybreaks
\usepackage{bookmark}
\IfFileExists{xurl.sty}{\usepackage{xurl}}{} % add URL line breaks if available
\urlstyle{same}
\hypersetup{
  pdftitle={计算量},
  colorlinks=true,
  linkcolor={Maroon},
  filecolor={Maroon},
  citecolor={Blue},
  urlcolor={Blue},
  pdfcreator={LaTeX via pandoc}}

\title{计算量}
\author{}
\date{}

\begin{document}
\maketitle

\subsubsection{7.2.3 计算量}\label{ux8ba1ux7b97ux91cf}

下面考虑解式（7.2.1）的 Gauss 消去法的计算量。

（1）消元过程的计算量。第一步计算乘数 \(m_{i1}\)（\(i=2,3,\cdots,n\)）需要 \(n-1\) 次除法运算，计算 \(a_{ij}^{(2)}\)（\(i,j=2,3,\cdots,n\)）需要 \((n-1)^2\) 次乘法运算及 \((n-1)^2\) 次加、减法运算。一般可列表（见表 7.1）计算。

\textbf{表 7.1}

\[
\begin{array}{c|c|c|c}
\text{第 }k\text{ 步}&\text{加、减法次数}&\text{乘法次数}&\text{除法次数}\\\hline
1&(n-1)^2&(n-1)^2&n-1\\
2&(n-2)^2&(n-2)^2&n-2\\
\vdots&\vdots&\vdots&\vdots\\
n-1&1&1&1\\\hline
\text{合计}&\dfrac{n(n-1)(2n-1)}{6}&\dfrac{n(n-1)(2n-1)}{6}&\dfrac{n(n-1)}{2}
\end{array}
\]

这里利用了求和公式

\[\sum_{i=1}^{n}i=n(n+1)/2,\quad \sum_{i=1}^{n}i^2=n(n+1)(2n+1)/6\quad(n\ge1).\]

消元过程所需的乘、除法次数 \(MD\) 及加、减法次数 \(AS\) 分别为

\[MD=n(n^2-1)/3,\quad AS=n(n-1)(2n-1)/6.\]

（2）计算 \(b^{(n)}\) 的计算量。

\[MD=(n-1)+(n-2)+\cdots+2+1=n(n-1)/2,\quad AS=n(n-1)/2.\]

（3）解 \(A^{(n)}x=b^{(n)}\) 所需的计算量。\(MD=n(n+1)/2\)，\(AS=n(n-1)/2\)。解式（7.2.1）所需的总的乘除法次数及加减法次数分别为

\[MD=n^3/3+n^2-n/3\approx n^3/3\quad\text{（当 $n$ 比较大时）},\]

\[AS=n(n-1)(2n+5)/6\approx n^3/3\quad\text{（当 $n$ 比较大时）}.\]

\textbf{定理 7.4} 如果 \(A\) 为 \(n\) 阶非奇异矩阵，则用 Gauss 消去法解式（7.2.1）所需乘除法次数及加减法次数分别为

\(1^\circ\;MD=n^3/3+n^2-n/3\)；

\(2^\circ\;AS=n(n-1)(2n+5)/6\)。

如果用 Cramer 法则解式（7.2.1），就需要计算 \(n+1\) 个 \(n\) 阶行列式，若行列式计

\href{../../source-images/pdf-184.jpeg}{原页扫描}

算是用子式展开，总共需要 \((n+1)!\) 次乘法运算。例如 \(n=10\) 时，Gauss 消去法需要 430 次乘、除法运算，而 Cramer 法则却需要 39 916 800 次乘法运算。由此可见，用 Cramer 法则解式（7.2.1）工作量太大，不便使用。

\end{document}
